\documentclass[10pt,a4paper]{amsart}
\usepackage[margin=19mm]{geometry}
\usepackage{amsmath,amssymb,mathtools}
\usepackage[scr=rsfs,cal=euler]{mathalfa}
\usepackage{xfrac}
\usepackage[unicode,hidelinks,bookmarks=false]{hyperref}
\newcommand{\ie}{i.e.\ }
\newcommand{\cf}{cf.\ }
\newcommand{\resp}{respectively\ }
\newcommand{\Subsection}{\subsection}
\providecommand{\nobreakdash}{\nobreak\hbox{-}}
\setlength{\emergencystretch}{2em}
\multlinegap0pt
\usepackage{bm}
\usepackage{bbm}
\usepackage{dsfont}
\usepackage{xspace}
\usepackage{cancel}
\usepackage{mathtools}
\usepackage{cleveref}
\usepackage[shortlabels]{enumitem}
\usepackage{amsthm}
\makeatletter
\@ifundefined{thm}{\newtheorem{thm}{Thm}}{}
\@ifundefined{definition}{\newtheorem{definition}[thm]{Definition}}{}
\@ifundefined{obs}{\newtheorem{obs}[thm]{Observation}}{}
\@ifundefined{prop}{\newtheorem{prop}[thm]{Proposition}}{}
\makeatother
\makeatletter
\expandafter\ifx\csname abstract\endcsname\relax
\renewenvironment{abstract}{%
      \if@twocolumn
        \section*{\abstractname}%
      \else
        \small
        \begin{center}%
          {\bfseries \abstractname\vspace{-.5em}\vspace{\z@}}%
        \end{center}%
        \quotation
        \noindent\ignorespaces
      \fi
    }{%
      \if@twocolumn\else\endquotation\fi
    }
\fi
\newcommand{\lrcornerqed}[0]{}
\makeatother
\title[Combinatorial enumeration of lattice paths by flaws with res]{Combinatorial enumeration of lattice paths by flaws with respect to a linear boundary of rational slope}
\author{Federico Firoozi, Jonathan Jedwab, Amarpreet Rattan}
\date{Source: arXiv:2406.09590v2 (CC BY 4.0)}
\begin{document}
\maketitle

\noindent\emph{Source and licence.} Combinatorial enumeration of lattice paths by flaws with respect to a linear boundary of rational slope. Version arXiv:2406.09590v2 (CC BY 4.0), distributed under the Creative Commons Attribution 4.0 International licence, as recorded in the arXiv licence metadata for this exact version and shown on the abstract page https://arxiv.org/abs/2406.09590v2. Full rights notice: licenses/arxiv-2406.09590-CC-BY-4.0.txt.\par\smallskip

\noindent\textit{Editorial scope.} Counted unit: the Prop whose source label is \texttt{rem:codomain\_psi\_X} in the article (source0.tex lines 2499--2502, source label \texttt{rem:codomain\_psi\_X}), delivered together with the article's own complete original proof of it (source0.tex lines 2503--2548, one \texttt{proof} environment of 46 lines). No mathematics is written, repaired or re-derived: statement and proof are the article's own text line for line. Required context retained uncounted: obs:flaws (lines 1306--1311); def:domain\_partition (lines 1803--1816); def:codomain\_partition (lines 1838--1851); def:codomain\_split (lines 1857--1875). Every definition, display and label of those lines is retained unchanged. Omitted from this package and not redistributed: 1--1305, 1312--1802, 1817--1837, 1852--1856, 1876--2498, 2549--2901. Nothing inside a retained excerpt is omitted or altered: every retained body is delivered exactly as the article's lines are written, so no omission receipt of any kind accompanies this package. Citations: the retained excerpts carry no citation command at all, so no bibliography entry of the article is transcribed and no reference is redistributed. Reference adaptation: none was needed and none was made; every reference inside the delivered unit resolves inside it, so no unresolved reference or citation is produced. Printed slips kept literal: no printed word, symbol, hypothesis or number of the counted statement or of its proof was altered. Layout and class: the article's own class is \texttt{article} with packages including etex, amsmath, amssymb, amsthm, mathtools, enumitem, tikz, colortbl, float, subcaption, booktabs, hyperref, cleveref, autonum; the wrapper is the lane's pinned \texttt{amsart} editorial wrapper at 10pt on A4, which restates the theorem-like environments the retained text needs and adds the article's own macro definitions that the retained text needs (\texttt{\textbackslash lrcornerqed}), with extra wrapper packages (bm, bbm, dsfont, xspace, cancel, mathtools, cleveref, enumitem). The delivered numbering is the unit's own. Assets: no retained span contains a figure environment, a graphics-inclusion command, a PDF-page inclusion, a table or a TikZ picture; no image file is copied into this package, the package declares no asset dependency and no image file is redistributed with it. Licence: version arXiv:2406.09590v2 of this article is distributed under the Creative Commons Attribution 4.0 International licence, as recorded in the arXiv licence metadata for this exact version and shown on the abstract page https://arxiv.org/abs/2406.09590v2. A full rights notice is delivered at licenses/arxiv-2406.09590-CC-BY-4.0.txt. Characters: every retained excerpt is pure ASCII, so the delivered engine, XeLaTeX with its default Latin Modern text font, reports no missing character.
\begin{obs}\label{obs:flaws}
    Let $p$ be a path containing exactly $\beta$ interior boundary points and exactly $\lambda$ LPAs.
    Suppose that $p$ is split at an HPB $H$ into $p_1p_2$, so that $H$ is the endpoint of $p_1$ and the startpoint of~$p_2$. Then the rearranged path $p_2p_1$ has exactly $\beta+1$ more flaws than~$p$, namely all $\beta$ interior boundary points of $p$ together with the endpoint of~$p_2$.
    If instead $p$ is split at an LPA into $p_1p_2$, then the rearranged path $p_2p_1$ has exactly $\lambda$ fewer flaws than~$p$, namely all $\lambda$ LPAs of~$p$.
    \lrcornerqed
\end{obs}

\begin{definition}[The subsets $X$ and $Y$]
\phantomsection
\label{def:domain_partition}
    Let $p \in W_k(g) \setminus S_k(g)$. Split $p$ at its last non-terminal boundary point into~$q r$, and regard $q$ and $r$ as paths.  
    The path $p$ lies in $Y$ provided:
    \begin{enumerate}[label=\textit{(\roman*)}]
        \item $q$ has at least one flaw, and
        \label{def:domain_partition:condition_1}
        
        \item the elevation of the LPAs of $q$ is smaller than the magnitude of the elevation of the HPBs of~$r$ (if any). 
        \label{def:domain_partition:condition_2}
    \end{enumerate}
    Otherwise, $p$ lies in~$X$. \lrcornerqed
\end{definition}

\begin{definition}[The subsets $\mathcal{X}$ and $\mathcal{Y}$]
\phantomsection
\label{def:codomain_partition}
    Let $\mathbbm{p} \in W_{k+1}(g)$. 
    The path $\mathbbm{p}$ lies in $\mathcal{Y}$ provided: 
    \begin{enumerate}[label=\textit{(\roman*)}]
        \item $\mathbbm{p}$ has at least two LPAs, and
        \label{def:codomain_partition:condition_1}
        
        \item the subpath of $\mathbbm p$ lying between the last two LPAs of $\mathbbm{p}$ contains no boundary points of~$\mathbbm{p}$.
        \label{def:codomain_partition:condition_2}
    \end{enumerate}
    Otherwise, $\mathbbm{p}$ lies in~$\mathcal{X}$. \lrcornerqed
\end{definition}

\begin{definition}[Canonical representation of paths in $\mathcal X$ and~$\mathcal Y$]
\phantomsection
\label{def:codomain_split}
    Let $\mathbbm{p} \in W_{k+1}(g)$.
    \begin{description}

        \item[Case $\mathbbm{p} \in \mathcal{X}$:]
            let $L$ be the last LPA of~$\mathbbm{p}$, and let $B$ be the boundary point of $\mathbbm{p}$ (possibly the startpoint of $\mathbbm{p}$) which immediately precedes~$L$. 
            Split $\mathbbm{p}$ at $B$ and $L$ into $\mathbbm{p} = \mathbbm{q} \mathbbm{r}_2 \mathbbm{r}_1$. 
            The canonical representation of $\mathbbm{p}$ is~$\mathbbm{q} \mathbbm{r}_2 \mathbbm{r}_1$.

        \item[Case $\mathbbm{p} \in \mathcal{Y}$:]
            the path $\mathbbm{p}$ has at least two LPAs. 
            Split $\mathbbm{p}$ at its 
            last two LPAs
            into $\mathbbm{p} = \mathbbm{q}_1 \mathbbm{r} \mathbbm{q}_2$. 
            The canonical representation of $\mathbbm{p}$ is~$\mathbbm{q}_1 \mathbbm{r} \mathbbm{q}_2$.\lrcornerqed
    \end{description}
\end{definition}

\begin{prop}
\label{rem:codomain_psi_X}
The map $\psi^\mathcal{X}$ is well-defined.
\end{prop}

\begin{proof}
    Let $\mathbbm{p} \in \mathcal{X}$.
    We must check that 
    $\psi^\mathcal X(\mathbbm{p}) = \mathbbm{q} \mathbbm{r}_1 \mathbbm{r}_2$ belongs to~$X$.

    Let $L$ be the endpoint of the path~$\mathbbm{r}_2$.
    By \cref{def:codomain_split}, we have:

    \begin{enumerate}
    \item[1.]
    $L$ is the last LPA of $\mathbbm{p}= \mathbbm{q} \mathbbm{r}_2 \mathbbm{r}_1$, and the startpoint of $\mathbbm{r}_2$ is the boundary point of $\mathbbm{p}$ which immediately precedes~$L$.
    \end{enumerate}
    
    By \cref{def:codomain_partition},  we have:

    \begin{enumerate}
    \item[2.]
    either $\mathbbm{p}$ has exactly one LPA,
    or the subpath of $\mathbbm{p}$ lying between the last two LPAs of $\mathbbm{p}$ contains a boundary point of~$\mathbbm{p}$.
    \end{enumerate}
    
    It follows from statements 1 and 2 that:

    \begin{enumerate}
    \item[3.]
    the subpath $\mathbbm{r}_2 \mathbbm{r}_1$ of $\mathbbm{p}$ contains exactly one LPA of $\mathbbm{p}$, namely the point~$L$.
    \end{enumerate}
    
    Statement 3 and \cref{obs:flaws} imply that:

    \begin{enumerate}
    \item[4.]
    the path $\psi^{\mathcal X}(\mathbbm{p}) = \mathbbm{q} \mathbbm{r}_1 \mathbbm{r}_2$ splits at its last non-terminal boundary point into the paths $\mathbbm{q}$ and~$\mathbbm{r}_1 \mathbbm{r}_2$.
    \item[5.]
    the path $\psi^\mathcal X(\mathbbm{p})= \mathbbm{q} \mathbbm{r}_1 \mathbbm{r}_2$ has exactly one fewer flaw than $\mathbbm{p}$.
    \end{enumerate}
    The elevation of $L$ in the path~$\mathbbm{r}_2 \mathbbm{r}_1$ equals the 
    magnitude of the elevation of the HPBs of the path~$\mathbbm{r}_1 \mathbbm{r}_2$. Since $L$ is an LPA of $\mathbbm{p}$ by statement 1, this gives:

    \begin{enumerate}
    \item[6.]
    the elevation of the LPAs of the path~$\mathbbm q$ (if any) is greater than or equal to the magnitude of the elevation of the HPBs of the path~$\mathbbm{r}_1 \mathbbm{r}_2$.
    \end{enumerate}
    
    Statements 4, 5 and 6 show by \cref{def:domain_partition} that $\psi^\mathcal X(\mathbbm{p}) \in X$.
\end{proof}

\end{document}
